% !TEX root = main2.tex
%\section{Other examples}
%\mitras{Created this section for describing the other examples. Could be merged with the modeling section or moved later.}

\subsection{ADAS and autonomous vehicle benchmarks}
\label{ssec:adas}
This is a suite of benchmarks we have created representing various common scenarios used for testing ADAS and Autonomous driving control systems. 
%
The hybrid system for a scenario is  constructed by putting together several individual vehicles.
The higher-level decisions (paths) followed by the vehicles are captured by transition graphs
while the detailed dynamics of each vehicle comes from a black-box 
\Simulink \ simulator from \Mathworks~\cite{simulinkcar}.

Each vehicle has several continuous variables including the $x,
y$-coordinates of the vehicle on the road, its velocity, heading, and
steering angle. The vehicle can be controlled by two input signals,
namely the throttle (acceleration or brake) and the steering speed.
%
By choosing appropriate values for these input signals, we have
defined the following modes for each vehicle ---
\Lmode{cruise}: move forward at constant speed, 
\Lmode{speedup}: constant acceleration,
\Lmode{brake}: constant (slow) deceleration,
\Lmode{em\_brake}: constant (hard) deceleration.
In addition, we have designed lane switching modes \Lmode{ch\_left} and \Lmode{ch\_right} in which the acceleration and steering are controlled in such a manner that the vehicle switches to its left (resp. right) lane in a certain amount of time. 

For each vehicle, we mainly analyze four variables: absolute position
($sx$) and velocity ($vx$) orthogonal to the road direction
($x$-axis), and absolute position ($sy$) and velocity ($vy$) along the
road direction ($y$-axis). The throttle and steering are captured using
the four variables. We will use subscripts to distinguish between
different vehicles.  The following scenarios are constructed by
defining appropriate sets of initial states and transitions graphs
labeled by the modes of two or more vehicles. In all of these
scenarios a primary safety requirement is that the vehicles maintain
safe separation. See Appendix~\ref{app:adas} for more details on
initial states and transition graphs of each scenario.
 %
\begin{description} 
\itemsep0em 
\item[$\auto{Merge}$:] 
Vehicle A in the left lane is behind   vehicle B  in the right lane. A switches through modes  \Lmode{cruise}, \Lmode{speedup}, \Lmode{ch\_right}, and \Lmode{cruise} over  specified intervals to merge behind B.
Variants of this scenario involve $B$ also switching to \Lmode{speedup} or \Lmode{brake}.
% within a time bound and maintains at least a given safe separation.
%%
%%%%
%%%%
%%
%\item[$\auto{MergeAhead}$:] 
%Initial condition: same as  $\auto{MergeBehind}$ with 
%except that B is in \Lmode{cruise} or \Lmode{brake} mode.
%Transition graph: same structure as  $\auto{MergeBehind}$ with different  timing parameters.
%
%Requirement: A merges ahead of B and maintains at least a given safe separation. 
%%
%%%%
%%%%
%%
\item[$\auto{AutoPassing}$:]
Vehicle A starts behind  B in the same lane, and  goes through a sequence of modes to  overtake B. If B switches to \Lmode{speedup} before A enters \Lmode{speedup} then A aborts and changes back to right lane. 
%Requirement: A maintains safe separation with  B.
\item[$\auto{Merge3}$:]
Same as $\auto{AutoPassing}$ with a third car C always ahead of $B$.
\item[$\auto{AEB}$:]
Vehicle A cruises behind B and B stops.
A transits from \Lmode{cruise} to \Lmode{em\_brake} possibly over several different time intervals as governed by different sensors and reaction times.

\end{description}



% source
% brief description
%% dynamics
%% graph
%% requirements

% results

% remarks

 
